
\minitoc 

% ========================================================================================================
% Introduction 

\subsection{Introduction}

Dans ce chapitre, nous nous intéresserons à la simulation de processus de Markov grâce 
à l'étude de celles-ci. Pour rappel, une chaîne de Markov décrit l'évolution d'une quantité soumise à des 
fluctuations aléatoires. C'est une notion plus générale que les variables iid. 

% ========================================================================================================
% Définition et propriétés 

\section{Définition et propriétés}

\begin{definition}[Chaîne de Markov]
    Une chaîne de Markov est une processus aléatoire $(X_n)_{n \geqslant 1}$ défini sur un espace $E$ 
    tel qu'il existe $ f \in \mathcal{M}(E \times [0,1], E)$ et $(u_n)_{n \geqslant 0}$ une suite idd de 
    loi $ \mathcal{U}([0,1])$ telle que : 
        $$ \boxed{\forall n \geqslant 1, \quad X_{n+1} = f(X_n, u_{n+1})} $$   
    avec : 
        $$ X_0 \perp (u_n)_{n \geqslant 0} $$
\end{definition}

\begin{example}
    Soit $P \in \mathcal{M}_n(\R)$ une matrice stochastique. Soit une chaîne de Markov sur 
    un espace $E = \llbracket 1, n \rrbracket $ de matrice de transition $P$. 
    Alors c'est bien une chaîne de Markov au sens de la définition précédente. 
    En effet, on pose : 
        \[ \forall i, j \in E, \quad f(i,x) = j \quad \text{avec} \quad x \in \left[ \sum_{k=1}^{j-1} P(i,j); \sum_{k=1}^{j} P(i,k)\right]  \] 
    On aura donc : 
        \begin{align*}
            & \myP(X_0 = x_0, \dots, X_n = x_n) \\
            &= \myP(X_0 = x_0, f(x_0, u_1) = x_1, \dots, f(x_{n-1}, u_n) = x_n) \\
            &= \myP(X_0 = x_0) \myP(f(x_0, u_1) = x_1) \dots \myP(f(x_{n-1}, u_n) = x_n) \\ 
            &= \myP(X_0 = x_0) P(x_0, x_1) \times \dots \times P(x_{n-1}, x_n)       
        \end{align*}
\end{example}

\begin{example}
    Soit le processus $(X_n)_{n \in \N}$ décrivant l'évolution de la température locale au cours du temps défini par : 
        $$ X_0 = 25 \quad X_{n+1} = \frac{X_n + 25}{2} + Z_{n+1} $$
    où $(Z_n)_{n \in \N}$ est une suite de variables aléatoires iid de loi $ \mathcal{N}(0, 10)$. 
    Alors ce processus est une chaîne de Markov avec : 
        $$ f(x,u) = \frac{x + 25}{2} + \Phi^{-1}(u) $$  
    avec $N \sim \mathcal{N}(0, 10)$ et : 
        \begin{align*}
            \Phi(u) = \myP(N \leqslant u) = \left( \int_{- \infty}^{u} e^{x^2 / 20} \; dx\right) \times \frac{1}{\sqrt{20 \pi}}   
        \end{align*}
\end{example}

\begin{remark}
    Cet exemple nous pose une problématique. En effet, nous avons définit notre chaîne de Markov avec seulement 
    une suite de variables iid de loi uniforme et chaque terme de la chaîne ne dépend que d'une seule 
    réalisation de cette loi uniforme. Pour nous permettre de faire dépendre un $X_n$ d'une loi normale, 
    nous devrions pouvoir utiliser deux variables uniforme pour la simuler (par la méthode de Box-Muller). 
\end{remark}

Le lemme suivant nous l'assure donc en affirmant que l'on peut produire autant de loi uniformes à partir d'une 
seule. 

\begin{lemma}[Doublement de la loi uniforme]
    Il existe $g : [0,1] \longrightarrow [0,1]^2$ mesurable telle que si $ U \sim \mathcal{U}([0,1])$ 
    et en posant $(V, W) = g(U)$ alors : 
        $$ V \perp W \quad \text{et} \quad W,V \sim \mathcal{U}([0,1]) $$  
\end{lemma}

\begin{corollary}[Génération de lois uniformes]
    Il existe une fonction mesurable $h : [0,1] \longrightarrow [0,1]^\N$ telle que $h(U)$ 
    est une suite de variables idd de loi $ \mathcal{U}([0,1])$.  
\end{corollary}

\begin{remark}
    Ces deux derniers résultats nous permettent d'écrire : 
        $$ X_{n+1} = f(X_n, Z_{n+1, 1}, \dots, Z_{n+1, k}) $$ 
    avec $(Z_{n,k})_{\substack{k \geqslant 1 \\ n \geqslant 1}}$ une matrice de suites iid. 
    On a alors que $(X_n)_{n \in \N}$ est une chaîne de Markov. 
\end{remark}

\begin{example}[Processus de Galton-Watson]
    Soit $(X_{n, k})_{\substack{n \geqslant 1, k \geqslant 1}}$ un tableau de variables 
    aléatoires iid à valeurs dans $\N$. Le processus défini par : 
        $$ Z_0 = 1 \quad \text{et} \quad Z_{n+1} = \sum_{j=1}^{Z_n} X_{n+1, j}  $$
    est une chaîne de Markov. C'est un exemple de processus représentant l'évolution d'une population 
    au cours du temps. Ici, $Z_n$ représente le nombre d'individus à la génération $n$ et $X_{n+1, j}$ le nombre 
    de descendants de $j$-ème individu de cette génération. 
    $Z_{n+1}$ est donc le nombre d'individus créés à la génération $n+1$. 
\end{example}

\paragraph{Notation. } Soit $(X_n)_{n \geqslant 1}$ une chaîne de Markov. Pour tout $x \in E$, on pose $ \myP_x$ la loi de 
$(X_n)_{n \geqslant 1}$ conditionnellement à $X_0 = x$. On a donc : 
    $$ 
        \begin{cases}
            X_0 = x \\ 
            X_{n+1} = f(X_n, u_{n+1})
        \end{cases}
    $$

\begin{prop}[Propriété de Markov]
    Soit  $(X_n)_{n \geqslant 1}$ une chaîne de Markov et $N \in \N$, alors : 
        $$ \mathbb{E}(h(X_{n+1}, X_{n+1}, \dots, X_{n+k}) \; | \; X_0 = x_0, \dots, X_N = x_N) 
        = \mathbb{E}_{X_n}(h(X_1, X_2, \dots, X_k)) $$   
\end{prop}

\begin{example}
    On définit par récurrence : 
        $$ F_0 = 1, F_1 = 1 \quad \text{et} \quad F_{n+1} = F_{n} + F_{n-1} + \varepsilon_{n+1} $$ 
    avec $( \varepsilon_n)_{n \in \N}$ une suite de variables aléatoires idd de loi $ \mathcal{U}(\{-1,0, 1\})$. 
    La suite $(F_n)_{n \in \N}$ représente la suite de Fibonacci pour laquelle on ferait des erreurs de calculs. 
    Cette suite n'est pas une chaîne de Markov car : 
        $$ \myP(F_4 = 2 \; | \; F_3 = 1) \not = \myP(F_4 = 2 \; | \; F_3 = 1, F_2 = 1) $$ 
    elle ne satisfait pas la propriété de Markov. Par contre, le couple $(F_n, F_{n-1})_{n \geqslant 1}$ est une 
    chaîne de Markov avec : 
        $$
            \begin{pmatrix}
                F_{n+1} \\ 
                F_n 
            \end{pmatrix}
            = 
            f \left(
                \begin{pmatrix}
                    F_n \\ 
                    F_{n_1}
                \end{pmatrix}, 
                \varepsilon_{n+1}
            \right)
        $$ 
    avec : 
        $$ f : 
            \begin{cases}
                \R^2 \times \R \longrightarrow \R^2 \\ 
                (x,y, \varepsilon) \longmapsto (x+y+ \varepsilon, x)
            \end{cases}
        $$
\end{example}

\begin{remark}
    Soit $(Z_n)_{n \geqslant 1}$ une suite de variables aléatoires idd et : 
        $$ S_n = Z_1 + \dots + Z_n $$ 
    alors : 
        \begin{itemize}
            \item $S_n$ est une chaîne de Markov 
            \item $ \max \{S_k, k \leqslant n\}$ n'est pas une chaîne de Markov. 
            \item $(S_n, \max \{S_k, k \leqslant n\})$ est une chaîne de Markov. 
        \end{itemize}
\end{remark}

\begin{proposition}[Algorithme de simulation]
    Comme dit précédement, une chaîne de Markov est définie telle que : 
        $$ 
            \begin{cases}
                X_0 = x_0 \\ 
                X_{n+1} = f(X_n, u_{n+1})
            \end{cases}    
        $$
    On peut donc utiliser l'algorithme suivant pour les simuler : 
\begin{lstlisting}
def simul_chaine_Markov(n,x0,f): 
    traj = [X0]
    for i in range(n): 
        traj.append(f(traj[-1],rand()))
    return traj
\end{lstlisting}
    La principale difficulté est donc de définir la fonction $f$. 
\end{proposition}

% ========================================================================================================
% Comportement asymptotique d'une chaîne de Markov 

\section{Comportement asymptotique d'une chaîne de Markov}

L'utilisation des chaînes de Markov en simulation est liée au théorème suivant. 

\begin{theorem}[Ergodique de Birkhoff]
    SOit $(X_n)_{n \geqslant 1}$ une chaîne de Markov. Si cette chaîne admet une mesure de probabilités $ \pi$ invariante 
    et ergodique alors $ \forall f in \in L^1( \pi)$, on a : 
        $$ \boxed{ \lim_{n \to \infty} \frac{1}{n} \sum_{i=1}^{n} f(X_i) = \int f \; d_{ \pi} \quad \text{p.s}  } $$
\end{theorem}

Par conséquent, pour calculer $ I = \int f \; d_\pi$ avec $ \pi$ une mesure coûteuse à simuler mais une mesure 
invariante et ergodique d'une chaîne de Markov, on peut approcher $I$ par $ \lim_{n \to \infty} \frac{1}{n} \sum_{i=1}^{n} f(X_i)$. 
Évidemment, puisque les variables $X_i$ ne sont pas iid, la convergence a lieu a un rythme plus lent. Sous certaines 
conditions, on peut avoir : 
    $$ \sqrt{n} \left( \frac{1}{n} \sum_{i=1}^{n} f(X_i) - \int f \; d_\pi  \right) \underset{n \to \infty}{\longrightarrow} \mathcal{N}(0,1) $$

\begin{definition}[Mesure Invariante]
    Soit $(X_n)_{n \in \N}$ une chaîne de Markov et $u \sim \mathcal{U}([0,1])$. 
    Soit $ \pi$ une mesure sur $(E, \mathcal{F}) $. On dit que $ \pi$ est \emph{invariante} si pour tout $A \in \mathcal{F}$ 
    on a : 
        $$ \int \myP_x(X_1 \in A) \; d_{ \pi(x)} = \pi(A) $$ 
        $$ \iff \int \myP(f(x,u_1) \in A) \; d_{ \pi(x)} = \pi(A) $$  
    Si $ \pi$ est une \emph{mesure de probabilité}, cette définition est équivalente à : 
        $$ X_0 \sim \pi \quad \text{et} \quad U_1 \sim \mathcal{U}([0,1]) \Longrightarrow X_1 = f(X_0, U_1) \sim \pi $$  
\end{definition}

\begin{example}
    Reprenons notre exemple de simulation de l'évolution de la température : 
        $$ X_{n+1} = \frac{X_n + 25}{2} + Z_n \quad \text{avec} \quad Z_n \sim \mathcal{N}(0,1) $$ 
    Si $X_0 \sim \mathcal{N}( \mu, \sigma^2)$ alors : 
        $$ X_1 = \frac{X_0 + 25}{2} + Z_1 \sim \mathcal{N} $$
    comme combinaison de gaussiennes indépendantes. De plus on a : 
        $$ \mathbb{E}(X_1) = \frac{25 + \mu}{2} \quad \text{et} \quad \V(X_1) = \frac{ \sigma^2}{4} + 10 $$    
    On a donc $ X_1 \sim \mathcal{N} \left(\frac{25 + \mu}{2}, \frac{ \sigma^2}{4} + 10\right)$. 
    La loi $ \mathcal{N}( \mu, \sigma^2)$ est donc invariante si : 
        $$ \mu = \frac{25 + \mu}{2} \quad \text{et} \quad \sigma^2 = \frac{ \sigma^2}{4} + 10 $$ 
        $$ \iff \mu = 25, \; \sigma^2 = \frac{40}{3} $$   
\end{example}

\begin{remark}
    Attention, si $X_0 \sim \pi$ avec $ \pi$ une mesure invariante, alors $ \forall n \in \N^*, X_n \sim \pi$ 
    mais $ X_n \not \perp X_{n-1}$ en général. 
\end{remark}

\begin{definition}[Mesure Ergodique]
    Une mesure $ \pi$ sur un espace mesurable $(E, \mathcal{F})$ est dite \emph{ergodique} pour une chaîne de Markov 
    $(X_n)_(n \in \N)$ si pour tout $A \in \mathcal{F}$ on a : 
        $$ \boxed{\forall x \in E, \quad \myP_x(X_1 \in A) = \mathbbm{1}_{x \in A} \Longrightarrow \pi(A) = 0 \; \text{ou} \; \pi(\overline{A}) = 0 }$$      
\end{definition}

Une mesure ergodique est une mesure sur $E$ qui ne peut pas "partager" sa masse entre différents sous-ensembles 
invariants pour la chaîne. Autrement dit, pour tout $A \subset E$ tel que, partant de $A$, la chaîne reste dans 
$A$ et partant de $\overline{A}$, la chaîne reste dans $\overline{A}$, la mesure ergodique charge soit 
$A$, soit $\overline{A}$, mais pas les deux.

Pour une chaîne de Markov dans $E$, cela implique que la chaîne ne peut pas rester bloquée entre plusieurs 
sous-ensembles invariants : elle explore tous les états accessibles dans l’ensemble invariant où elle a démarré, 
et finit par "oublier" son point de départ.

\begin{remark}
    Si une chaîne de Markov n'est pas ergodique, cela signifie qu'il existe 
    au moins un sous-ensemble invariant non trivial $A \subset E$ tel que : 
    \begin{itemize}
        \item partant de $A$, la chaîne reste dans $A$ ;
        \item partant de $A^c$, la chaîne reste dans $A^c$.
    \end{itemize}

    Une mesure invariante peut alors mettre du poids à la fois sur $A$ et sur $A^c$. 
    En pratique, cela signifie que la chaîne peut rester confinée dans différents 
    "blocs" selon l'état initial, et donc qu'elle ne se mélange pas complètement.  

    Autrement dit, une chaîne non ergodique peut être décomposée en plusieurs 
    sous-chaînes indépendantes (les ensembles invariants), et la dynamique globale 
    dépend de l'état initial.
\end{remark}

\begin{prop}[Ergodicité et Irréductibilité]
    Pour les chaînes de Markov sur des espaces d'états dicrets, être ergodique est équivalent à être irréductible. 
\end{prop}

\begin{corollary}[Sur les espaces d'états dicrets]
    Si $E$ est un espace d'états dicret, alors toute chaîne de Markov irréductible est $ \mu$-ergodique pour toute 
    mesure de probabilité $ \mu$. 
\end{corollary}

\begin{theorem}[Birkhoff (Variante)]
    Soit $(X_n)_{n \in \N}$ une chaîne de Markov de mesure invariante et ergodique $ \pi$. 
    Soient $f,g \in L^1( \pi)$ on a alors : 
        $$ \boxed{ \lim_{n \to \infty} \frac{\sum_{i=1}^{n} f(X_i)}{\sum_{i=1}^{n} g(X_i)} = \frac{\int f \; d \pi}{\int g \; d \pi} \quad \text{p.s} } $$
\end{theorem}

Intuitivement, puisque qu'une chaîne ergodique se "mélange bien" dans son espace d'état, alors 
les fréquences empiriques (la moyenne des trajectoire) reflètent les fréquences théoriques (la mesure $ \pi$). 

\begin{remark}
    Ce théorème, en plus d'être une monstre ambulant de théorie des chaînes de Markov est très fort. 
    En effet, on obtient un résultats similaire à la loi des grands nombres alors que les $X_i$ ne sont pas 
    indépendants. Ici, même si $ \pi$ est difficile à calculer, on peut estimer $\int f \; d \pi$ 
    juste en simulant une trajectoire de la chaîne et en en faisant la moyenne. 
    Enfin, $f$ et $g$ peuvent être des fonctions assez quelconques tant qu'elles sont $L^1$. 
\end{remark}

\begin{example}
    Soit $(S_n)_{n \geqslant 1}$ la marche aléatoire simple sur $\Z$. Cette chaîne de Markov n'admet pas de 
    mesure de probabilité invariante. En revanche, la mesure de comptage sur $\Z$ est invariante et ergogique. Par conséquent : 
        $$ \forall A, B \subset \Z, \quad \lim_{n \to \infty} \frac{\sum_{k=1}^{n} \mathbbm{1}_{S_k} \in A}{\sum_{k=1}^{n} \mathbbm{1}_{S_k} \in B  } = \frac{|A|}{|B|} $$
    alors que : 
        $$ \frac{1}{n} \sum \mathbbm{1}_{S_k} \in A = 0 $$   
\end{example}


